\clearpage\begingroup\SourceLeaf{356b}\ArithmeticHead{25}{R}{DIE GRUNDLAGEN DER GEOMETRIE.}
\begin{center}\textbf{5.}\par\smallskip
\textls[100]{Schwanken und Zweifel (1799—1805).}
\end{center}
\noindent\hspace*{\parindent}Wir kehren zu den Beziehungen zwischen \textsc{Gauss} und \textsc{Bolyai} zurück. Im
Sommer 1799 nach Siebenbürgen zurückgekehrt, war \textsc{Wolfgang} zunächst
durch andere Geschäfte in Anspruch genommen worden und hatte die Mathematik
liegen lassen. Erst nachdem er im Frühjahr 1804 die Professur für
Mathematik und Physik am evangelisch-reformierten Kollegium zu Maros-Vásárhely
angetreten hatte, nahm er die »Göttingische Parallelentheorie« wieder
vor, feilte sie aus und sandte den Entwurf am 16. September 1804 an \textsc{Gauss}.
»Ich kann den Fehler nicht entdecken, prüfe Du der Wahrheit getreu und
\par\vfill\noindent{\footnotesize X\textsubscript{2} Abh. 4.\hfill 4\hspace*{2em}}
\clearpage\SourceLeaf{357}\ArithmeticHead{26}{L}{STÄCKEL, GAUSS ALS GEOMETER.}
\noindent schreibe mir so bald als nur möglich.\ .\ .\ .\ Wenn Du dieses Werkchen davor
wert hieltest (ich setze den Fall), so schicke es einer würdigen Akademie hin,
dass es beurteilt werde« (Br. G.-\textsc{Bolyai}, S. 65).
\par
Wie stellte sich \textsc{Gauss} zu den Bemühungen seines Freundes, das Parallelenaxiom
zu beweisen? Anders, als man es nach seinem Briefe vom 16. Dezember
1799 erwarten durfte. Er schreibt am 25. November 1804: »Ich habe
Deinen Aufsatz mit grossem Interesse und Aufmerksamkeit durchgelesen und
mich recht an dem echten gründlichen Scharfsinne ergötzt. Du willst aber
nicht mein leeres Lob, das auch gewissermassen schon darum parteiisch
scheinen könnte, weil Dein Ideengang sehr viel mit dem meinigen Ähnliches
hat, worauf ich ehemals die Lösung dieses Gordischen Knotens versuchte und
vergebens bis jetzt versuchte. Du willst nur mein aufrichtiges, unverhohlenes
Urteil. Und dies ist, dass Dein Verfahren mir noch \textls[100]{nicht} Genüge leistet.
Ich will versuchen, den Stein des Anstosses, den ich noch darin finde (und
der auch wieder zu derselben \textls[100]{Gruppe} von Klippen gehört, woran meine Versuche
bis jetzt scheiterten) mit so vieler Klarheit, als mir möglich ist, ans
Licht zu ziehen. Ich habe zwar noch immer die Hoffnung, dass jene Klippen
einst, und noch vor meinem Ende eine Durchfahrt erlauben werden. Indess
habe ich jetzt so manche andere Beschäftigungen vor der Hand, dass ich
gegenwärtig daran nicht denken kann, und glaube mir, es soll mich herzlich
freuen, wenn Du mir zuvorkommst und es Dir gelingt, alle Hindernisse zu
übersteigen. Ich würde dann mit der innigsten Freude alles tun, um Dein
Verdienst gelten zu machen und ins Licht zu stellen, so viel in meinen
Kräften steht« (W. VIII, S. 160).
\par
\textsc{Bolyai} hat diese Äusserungen als eine Ermunterung aufgefasst, sich weiter
um den Beweis zu bemühen, und mit Recht. »Meine Ideen gefielen ihm
überhaupt gar sehr, und er machte mich [in dem Brief vom 25. November
1804] darauf aufmerksam, welch hochwichtige Sache die Materie der Parallelen
sei, obwohl er davon [von der Göttingischen Parallelentheorie] doch keineswegs
befriedigt war« (\textsc{Bol.} S. 90). Am 27. Dezember 1808 sandte er an \textsc{Gauss}
einen Nachtrag (Br. G.-\textsc{Bolyai}, S. 96, vgl. \textsc{Bol.} S. 223). Als dieser keine
Antwort gab, ist der Briefwechsel bis zum Jahre 1816 unterbrochen worden.
Etwa bis zu diesem Jahre hat \textsc{Wolfgang} hart mit dem zweitausendjährigen
Problem gerungen, und hat schliesslich nichts davon getragen, als die Einsicht,
\clearpage\SourceLeaf{358}\ArithmeticHead{27}{R}{DIE GRUNDLAGEN DER GEOMETRIE.}
\noindent dass er alle seine Mühe verschwendet habe. »Schauderhafte, riesige
Arbeiten habe ich vollbracht, habe bei Weitem Besseres geleistet, als bisher
[geleistet wurde], aber keine vollkommene Befriedigung habe ich je gefunden;
hier aber gilt es: si paullum a summo discessit, vergit ad imum« (\textsc{Bol.} S. 77).
\par
Als \textsc{Johann Bolyai} von Wien aus, wo er seit 1818 Schüler der militärischen
Ingenieur-Akademie war, im Frühjahr 1820 dem Vater mitteilte, dass
er versuche, das elfte Axiom zu beweisen, war dieser aufs äusserste erschrocken
und beschwor ihn \textit{mit} den beweglichsten Worten, die Lehre von den Parallelen
in Frieden zu lassen. »Verliere keine Stunde damit. Keinen Lohn
bringt es, und es vergiftet das ganze Leben. Selbst durch das Jahrhunderte
dauernde Kopfzerbrechen von hundert grossen Geometern ist es schlechterdings
unmöglich, [das elfte] ohne ein neues Axiom zu beweisen. Ich glaube
doch alle erdenklichen Ideen diesfalls erschöpft zu haben. Hätte \textsc{Gauss} auch
fernerhin seine Zeit mit Grübeleien über dem elften Axiom zugebracht, so
wären seine Lehren von den Vielecken, seine \textit{Theoria motus corporum coelestium}
und alle seine sonstigen Arbeiten nicht zum Vorschein gekommen, und
er ganz zurückgeblieben. Ich kann es schriftlich nachweisen, dass er seinen
Kopf über die Parallelen zerbrach. Er äusserte mündlich und schriftlich, dass
er fruchtlos darüber nachgedacht habe« (\textsc{Bol.} S. 90).
\par
Durch die nachdrücklichen Warnungen seines Vaters wurde \textsc{Johann} nicht
abgeschreckt, im Gegenteil, seine Begierde, um jeden Preis durchzudringen,
wuchs auf das heftigste (\textsc{Bol.} S. 79). Gegen Ende des Jahres 1823 gelang es
ihm, den Gordischen Knoten zu durchhauen. Die unerwartete Lösung, die
er fand, war damals bereits im Besitz von \textsc{Gauss}, der, wie wir sehen werden,
nach langen Zweifeln um das Jahr 1816 zur Gewissheit gekommen war.
\par
Am Eingang des Briefes vom 25. November 1804 hatte \textsc{Gauss} angedeutet,
dass sein Ideengang Ähnlichkeit mit dem \textsc{Wolfgang}s habe. Dieser hatte die
Linie betrachtet, die entsteht, wenn man in gleichweit von einander abstehenden
Punkten einer Geraden nach derselben Seite Lote derselben Länge
errichtet und die auf einander folgenden Endpunkte durch Gerade verbindet.
Während man in der euklidischen Geometrie auf solche Art eine Parallele
zur Grundlinie erhält, ergibt sich in der nichteuklidischen Geometrie ein gebrochener
Linienzug, der aus gleich langen, unter gleichen Winkeln an einander
stossenden Strecken besteht. \textsc{Bolyai} hatte zu zeigen versucht, dass ein
\par\vfill\hfill{\footnotesize 4\textsuperscript{*}}\par
\clearpage\SourceLeaf{359}\ArithmeticHead{28}{L}{STÄCKEL, GAUSS ALS GEOMETER.}
\noindent Linienzug der angegebenen Art, wenn man weit genug auf ihm fortgehe, die
Grundlinie schneiden müsse; damit wäre nachgewiesen, dass die Annahme,
das elfte Axiom gelte nicht, auf einen Widerspruch führt. Dass \textsc{Gauss} sich
nach derselben Richtung hin versucht hat, wird durch eine Bemerkung bezeugt,
die sich auf der letzten Seite des Handbuches »\textit{Mathematische Brouillons}«
findet (W. VIII, S. 163); allerdings beginnen die Aufzeichnungen des Handbuches
erst mit dem Oktober 1805.
\par
In der Zeit zwischen 1799 und 1804 hatte \textsc{Gauss} aber noch auf einem
anderen Wege vorzudringen versucht. Notizen aus dem Jahre 1803 (W. X\textsubscript{1},
S. 451) geben mehrere Ansätze, mittels geometrischer Konstruktionen und
daraus abgeleiteter Funktionalgleichungen, also durch dasselbe Verfahren, das
\textsc{Gauss} auch bei dem Dreiecksinhalt angewandt hat (siehe S. 45), die zwischen
den Stücken eines Dreiecks geltenden Beziehungen herzuleiten. Damals sind
seine Anstrengungen vergeblich gewesen; vielleicht liegt hierin der Grund,
warum er gegenüber dem 1799 geäusserten Zweifel an der Wahrheit der Geometrie
im Jahre 1804 von der Hoffnung spricht, nach vor seinem Ende eine
Durchfahrt nach dem Hafen des Beweises für das Parallelenaxiom zu finden.
\clearpage\endgroup
